\documentclass[11pt]{article}
\usepackage{preamble}
\renewcommand{\answer}[1]{}

\begin{document}

\ladderheading{AP Statistics \textperiodcentered\ Unit 4 \textperiodcentered\ Topic 4.8 \textperiodcentered\ B: 2027-01-26 \textperiodcentered\ E: 2027-03-15}{A Controller Comparison}

\focusbanner{TODAY'S PROBLEM}

Suppose a food cooperative has 1{,}200 refrigerated display cases with conventional controllers and 900 with adaptive controllers that adjust cooling automatically. It takes independent random samples of 40 cases from each group; controller type was not assigned. A 95\% interval for conventional minus adaptive mean daily use is $(1.219,4.781)$ kWh per case. Replacing controllers is worthwhile only if mean savings are at least 2.0 kWh.\par\smallskip \textbf{Rhea:} ``Zero is excluded and 2 is inside, so replacing controllers is shown to cause savings of at least 2 kWh.''\par \textbf{Sol:} ``I would compare the interval with both 0 and 2, then check how the groups were formed.''\par
\begin{center}
\textbf{Daily kWh samples}\par\smallskip
\begin{tabular}{|l|c|c|c|c|}
\hline
\textbf{Controller} & \textbf{$N$} & \textbf{$n$} & \textbf{$\bar x$} & \textbf{$s$} \\ \hline
\textbf{Conv.} & 1,200 & 40 & 18.6 & 4.0 \\ \hline
\textbf{Adaptive} & 900 & 40 & 15.6 & 4.0 \\ \hline
\end{tabular}
\end{center}
\begin{firsttakebox}{FIRST TAKE --- before the video (5 min)}
Does this range seem to promise savings of at least 2 kWh? Explain in one sentence.
\workspace{0.9in}
\end{firsttakebox}

\begin{callout}{\IconBook\ CHECK ZERO, THE CLAIM'S THRESHOLD, AND THE DESIGN (keep on the page)}{calloutgreen}
A confidence interval for $\mu_1-\mu_2$ describes plausible values for a difference in
population means in the stated subtraction order, not a range containing a percent of individual
observations. If 0 is excluded, the interval supports a population mean difference; a practical
claim such as $\mu_1-\mu_2\geq c$ requires the entire interval to lie at or beyond $c$.
For a fixed interval, say ``we are confident,'' not ``there is a probability.'' Random sampling
supports population scope, while random assignment is needed to isolate a causal treatment effect.
\end{callout}

\newpage
\textbf{Finish after the video:}\par\smallskip

\dokbadge{1}~\textbf{(a)} Locate 0 and 2 relative to the interval.\par
\workspace{0.7in}

\dokbadge{2}~\textbf{(b)} Interpret the supplied 95\% interval in one sentence about the two populations of display cases.\par
\workspace{1.4in}

\dokbadge{3}~\textbf{(c)} In 3--4 sentences, decide whose account is justified. Use 0 and the 2-kWh cutoff, then explain why the way controllers were chosen limits a cause-and-effect claim.\par
\workspace{2.0in}

\begin{sentenceframebox}\raggedright\textbf{Frame:} Relative to the interval, 0 is \blank[0.9in] and 2 is \blank[0.9in]; those positions support \blank[2.4in].\end{sentenceframebox}

\begin{sentenceframebox}\raggedright\textbf{Frame:} The sampling design supports \blank[2.0in] but leaves \blank[1.8in] unresolved because \blank[1.8in].\end{sentenceframebox}

\begin{summaryexitbox}{EXIT --- turn this in}
One thing the video changed about my first take: \hrulefill\par
\workspace{0.4in}
\end{summaryexitbox}

\end{document}
